\begin{figure*}[t]
\centering
\begin{minipage}{.46\linewidth}
\centering
\begin{tikzpicture}[
  x=.64cm,y=12cm,
  accepted/.style={circle,draw=editgreen,fill=white,line width=1.15pt,minimum size=5.4pt,inner sep=0pt},
  rejected/.style={cross out,draw=fixgrey,line width=1.05pt,minimum size=5.4pt,inner sep=0pt},
  every node/.style={font=\scriptsize}
]
  \draw[black!18] (0,0) grid[xstep=1,ystep=.05] (10,.35);
  \draw[->,black!70] (0,0) -- (10.35,0) node[right] {Round};
  \draw[->,black!70] (0,0) -- (0,.39);
  \node[rotate=90] at (-.95,.18) {Validation F1};
  \foreach \x in {0,...,10} \draw (\x,0) -- +(0,-.005) node[below=2pt] {\x};
  \foreach \y/\lab in {0/.55,.05/.60,.10/.65,.15/.70,.20/.75,.25/.80,.30/.85,.35/.90}
    \draw (0,\y) -- +(-.04,0) node[left=2pt] {\lab};

  \fill[editfill,opacity=.55] (0,.36) rectangle (5.45,.385);
  \fill[blue!7] (5.55,.36) rectangle (10,.385);
  \node[font=\scriptsize\bfseries,text=editgreen] at (2.7,.3725)
    {rounds 1--5: $\tau_{\mathrm{imp}}=0$};
  \node[font=\scriptsize\bfseries,text=blue!55!black] at (7.8,.3725)
    {rounds 6--10: $\tau_{\mathrm{imp}}=.01$};
  \draw[densely dashed,red!65,line width=.9pt] (5.5,0) -- (5.5,.36);

  % Reported incumbent lineage.
  \draw[editgreen,line width=1.45pt]
    (0,.0229) -- (1,.0229) -- (1,.1849) -- (2,.1849) -- (2,.1855)
    -- (4,.1855) -- (4,.2413) -- (5,.2413) -- (5,.2823)
    -- (7,.2823) -- (7,.3005) -- (10,.3005) -- (10,.2591);

  \node[accepted] at (0,.0229) {};
  \node[accepted] at (1,.1849) {};
  \node[accepted] at (2,.1855) {};
  \node[accepted] at (4,.2413) {};
  \node[accepted] at (5,.2823) {};
  \node[accepted] at (7,.3005) {};
  \node[accepted] at (10,.2591) {};

  % Rejections use the compact baseline row convention of the Self-Harness
  % trajectory figure; their measured scores and context are in the side panel.
  \foreach \x/\lab in {3/R3,6/R6,8/R8,9/R9} {
    \node[rejected] at (\x,.010) {};
    \node[text=fixgrey,anchor=south] at (\x,.017) {\lab};
  }

  % Separately executed continuations, excluded from the reported lineage.
  \draw[fixgrey!62,densely dashed,line width=1.05pt]
    (2,.1855) -- (3,.2338) -- (4,.2338);
  \node[circle,draw=fixgrey!75,fill=white,line width=1pt,minimum size=5pt,inner sep=0pt] at (3,.2338) {};
  \node[rejected] at (4,.2338) {};

  \draw[fixgrey!62,densely dashed,line width=1.05pt]
    (5,.2823) -- (6,.2868) -- (7,.2442);
  \node[circle,draw=fixgrey!75,fill=white,line width=1pt,minimum size=5pt,inner sep=0pt] at (6,.2868) {};
  \node[rejected] at (7,.2442) {};

  % Marker legend, placed below the iteration axis as in Self-Harness.
  \node[accepted] at (1.30,-.060) {};
  \node[anchor=west] at (1.43,-.060) {accepted};
  \node[rejected] at (3.90,-.060) {};
  \node[anchor=west] at (4.03,-.060) {rejected};
  \draw[fixgrey!62,densely dashed,line width=1.05pt] (6.05,-.060) -- (6.45,-.060);
  \node[anchor=west] at (6.55,-.060) {abandoned branch};
\end{tikzpicture}
\end{minipage}%
\hfill
\begin{minipage}{.52\linewidth}
\centering\scriptsize
\setlength{\tabcolsep}{2.2pt}
\renewcommand{\arraystretch}{1.05}
\begin{tabularx}{\linewidth}{@{}lrrrrrX@{}}
\toprule
R & F1 & $\Delta$F1 & $\Delta$tok & Fail. & Dec. & What changed \\
\midrule
1  & .7349 & +.1620 & $-8.9\%$   & $23\to6$ & Prom. & emit user-ID pairs \\
2  & .7355 & +.0007 & $+11.7\%$  & $6\to1$  & Prom. & recover missing users \\
3  & .7838 & +.0483 & $-3.4\%$   & $1\to1$  & Rej.  & wrong-format regression \\
4  & .7913 & +.0558 & $-12.1\%$  & $1\to0$  & Prom. & confirm target categories \\
5  & .8323 & +.0410 & $+16.8\%$  & $0\to0$  & Prom. & correct detected errors \\
6  & .8368 & +.0045 & $-14.1\%$  & $0\to0$  & Rej.  & below revised $\tau_{\mathrm{imp}}=.01$ \\
7  & .8505 & +.0182 & $-33.8\%$  & $0\to0$  & Prom. & resolve count mismatch \\
8  & .8164 & $-$.0341 & $+56.5\%$  & $0\to1$  & Rej.  & format regression, +1 fail \\
9  & .6546 & $-$.1959 & $+665.7\%$ & $0\to7$  & Rej.  & token blowup, +7 fail \\
10 & .8091 & +.0128 & $-11.0\%$  & $0\to0$  & Prom. & rebuild after corrections \\
\bottomrule
\end{tabularx}
\end{minipage}
\caption{OOLONG-Pairs optimization trajectory (left) and the full per-round
record it summarizes (right). Changes are relative to that round's incumbent;
terminal failure counts are out of 72 attempts. Dark-green circles and the
step line denote accepted candidates on the reported incumbent lineage; gray
crosses in the row above the iteration axis denote rejected candidates.
Light-gray dashed traces show two abandoned exploratory continuations, found
by checking every sibling experiment directory in this lineage.}
\label{fig:optimization-trajectory}
\end{figure*}

% Diff-card primitives for fig:final-surface-edits: a full-width shaded row
% per logical line, "+" for appended text in that round's own color, plain
% (phantom-indented to match) for text retained verbatim from H0.
% Colorbox-around-parbox (not soul's \hl) so the shading spans the whole
% card width even when a line wraps. Round-to-text attribution verified
% against the real harness.json snapshots at each round (not the short
% behavioral_change summaries, which omit one of round 10's two insertions).
\newcommand{\diffsame}[1]{%
  \noindent\parbox{\dimexpr\linewidth-2\fboxsep\relax}{\ttfamily\small\phantom{+ }#1}\par\vspace{2pt}}
\newcommand{\diffadd}[2]{%
  \noindent\colorbox{#1}{\parbox{\dimexpr\linewidth-2\fboxsep\relax}{\ttfamily\small+\ #2}}\par\vspace{2pt}}
% Round-number badge for card titles: black text, background matching that
% round's own row fill color (not colored text) -- tight padding so it reads
% as a small tag rather than a full-size colorbox.
\newcommand{\rbadge}[2]{{\setlength{\fboxsep}{1.5pt}\colorbox{#1}{\textcolor{black}{#2}}}}

\begin{figure*}[p]
\centering
\setlength{\fboxsep}{4pt}
\begin{tcolorbox}[
  enhanced, colback=white, colframe=black!35, boxrule=0.5pt,
  arc=2pt, left=3pt, right=3pt, top=3pt, bottom=4pt,
  colbacktitle=black!8, coltitle=black, fonttitle=\bfseries\small,
  title={S2---decomposition instruction (promoted round \rbadge{r4fill}{4})}
]
\diffsame{Start by probing `context' to understand what you have, then decide how the task breaks into steps.}
\diffadd{r4fill}{Before selecting target categories, re-read the task question from `context' to identify which categories are required.}
\diffadd{r4fill}{Store the required categories in a variable and print them to confirm they match the task.}
\end{tcolorbox}

\vspace{4pt}
\begin{tcolorbox}[
  enhanced, colback=white, colframe=black!35, boxrule=0.5pt,
  arc=2pt, left=3pt, right=3pt, top=3pt, bottom=4pt,
  colbacktitle=black!8, coltitle=black, fonttitle=\bfseries\small,
  title={S3---execution instruction (cumulative promoted edits from rounds
    \rbadge{r1fill}{1}, \rbadge{r2fill}{2}, \rbadge{r5fill}{5}, and
    \rbadge{r10fill}{10})}
]
\diffsame{On each turn, run one block of code, print what you need to see, and build the result up in REPL variables.}
\diffadd{r1fill}{After classifying questions, generate the required user-ID pairs and output them in the format (user\_id\_1, user\_id\_2), one per line.}
\diffadd{r10fill}{When generating pairs from two user sets, skip pairs where the two user IDs are equal (a user cannot pair with itself).}
\diffadd{r2fill}{After classifying, verify that all users with entity or human being questions are captured. If any user appears to be missing, reclassify their questions.}
\diffadd{r5fill}{If you detect a classification error (e.g., a user's question count doesn't match the expected pattern), reclassify the affected questions before proceeding.}
\diffadd{r10fill}{After any reclassification, you MUST rebuild the user-to-category mapping and regenerate all pairs from scratch using the updated classifications. Do not reuse a previously computed user set or pair list after any classification change.}
\end{tcolorbox}

\vspace{4pt}
\begin{tcolorbox}[
  enhanced, colback=white, colframe=black!35, boxrule=0.5pt,
  arc=2pt, left=3pt, right=3pt, top=3pt, bottom=4pt,
  colbacktitle=black!8, coltitle=black, fonttitle=\bfseries\small,
  title={S4---verification instruction (promoted round \rbadge{r7fill}{7})}
]
\diffsame{Check your candidate answer before setting `answer["ready"] = True'.}
\diffadd{r7fill}{If you compute an expected count (e.g., of pairs, users, or classifications) and it does not match the actual count, do not submit until the discrepancy is resolved.}
\diffadd{r7fill}{Print the source of the discrepancy: which users or classifications differ between the expected and actual sets.}
\diffadd{r7fill}{Correct the underlying issue (e.g., reclassify affected questions, regenerate pairs) and re-verify before submitting.}
\end{tcolorbox}

\caption{Full retained field values for every surface changed in the frozen
round-10 harness, shown as a diff against the initial $H_0$ text: a plain line
is retained verbatim from $H_0$; a ``+'' line is new text, colored by the
round that added it (round number colored to match in each card's title).
S3 draws on four rounds, each shown in its own color; round 10 contributes
two separate insertions, one in the middle of the existing text and one at
the end, not a single trailing append. Every edit to these three surfaces was
a pure addition -- no span of the original text was altered or removed, so
no ``-'' lines appear anywhere. Unlike Figure~\ref{fig:optimization-trajectory},
this figure does not summarize the edits: it reproduces the full text from
the saved harness artifacts, diffed directly against the real per-round
harness snapshots (not the short per-candidate change summaries, which omit
one of round 10's two insertions).}
\label{fig:final-surface-edits}
\end{figure*}
